% RESEARCH_DRAFT_SHA256: 9d87d068fc71b8940e496675449095a9a6bc0e61d80fa54436f951466a22d1f3
% Layout and source-reference wrappers expanded for the web renderer.
\documentclass[11pt]{article}
\usepackage[margin=25mm]{geometry}
\usepackage{fontspec,unicode-math}
\setmainfont{Latin Modern Roman}
\setmathfont{Latin Modern Math}
\usepackage{amsmath,amssymb,mathtools,xurl,hyperref}
\hypersetup{colorlinks=true,urlcolor=blue,linkcolor=blue}
\setcounter{secnumdepth}{0}
\setlength{\parindent}{0pt}
\setlength{\parskip}{5pt}
\setlength{\emergencystretch}{3em}
\newcommand{\sref}[2]{\hyperlink{source:#1:#2}{[S#2]}}
\newcommand{\cataloguescope}[1]{\paragraph{Source record.}#1}
\newcommand{\floor}[1]{\left\lfloor#1\right\rfloor}
\begin{document}
\section*{Research-attempt provenance}
This standalone research attempt retains the complete original definition body below,
including its historical source and preparation notes. Those notes describe the
earlier catalogue preparation; the new attempt and its verification are identified
separately after that body. No original ID, catalogue status, or source statement is changed.

% UnsolvedMath ID 20002151; source catalogue code AIM-INFRASTRUCTURE-0051
\section{Context-closed fragments and version-pinned overlays}\label{20002151}
{\small UnsolvedMath ID 20002151 · Computer Science}
\subsection{Definitions and mathematical statement}

\paragraph{Shared notation.}Standard number systems are $\mathbb N,\mathbb Z,\mathbb Q,\mathbb R,\mathbb C$. Specialized notation and index ranges follow the cited source and the definitions below.


A document fragment is a selected subtree or portion of a larger source. A preamble supplies declarations and configuration used by a document processor. Standalone compilation supplies the context needed to process a fragment independently; inclusion processes it within the enclosing context. An overlay journal adds editorial material around papers hosted elsewhere.
\paragraph{Record scope and qualifications.}
The catalogue title is retained as metadata. Named certificates, contracts, bounds or frameworks in that title are not asserted as results of the original source extract.
The fragment boundaries, configuration inheritance, versions, and overlay-journal requirements are left open. This is a document-engineering proposal.

Catalogue screening limitation: PreTeXt software design, documentation, or publication feature request; no mathematical problem is specified by the supplied passage.

The source snapshot labels this record partially solved; this is catalogue metadata rather than independent verification of a mathematical result or the current completion of a project.
\paragraph{Statement recorded in the supplied source.}
"Standalone" as a way to compile a fragment, and ignore the preamble when in a large document. (Can this be managed by "versions"?). Related, an Overlay journal: material wrapped around a paper that appears elsewhere. \href{https://github.com/PreTeXtBook/community/wiki/Moderated-Problem-Session}{[S2]}
\paragraph{Source status.}
The supplied catalogue labels this source record partially solved; that metadata was checked against the supplied snapshot on 2026-10-11. The label is preserved for provenance and is not an independent verification of present open status or mathematical correctness.
\subsection{Short English statement}
Design standalone fragment compilation and consider related publication wrappers.
\subsection{Sources}
\begin{itemize}
\item[S1] Ulam AI and the UnsolvedMath Contributors, UnsolvedMath: A Curated Collection of Open Mathematics Problems, record 20002151 (AIM-INFRASTRUCTURE-0051). Catalogue curation and metadata: CC BY 4.0. The source extract is retained; definitions, English summaries and qualifications are editorial additions.
\url{https://huggingface.co/datasets/ulamai/UnsolvedMath}.
\item[S2] AIM Workshop Problem Lists, source linked by catalogue record AIM-INFRASTRUCTURE-0051.
\url{https://github.com/PreTeXtBook/community/wiki/Moderated-Problem-Session}.
\end{itemize}
\paragraph{Preparation and review.}
The original source and UnsolvedMath contributors retain attribution for the source extract. Codex (GPT-6; exact serving revision not exposed) assisted with definitions, English summaries, scope notes and formatting. No new human mathematical authorship or independently verified solution is asserted. Permanent source ID 20002151 and code AIM-INFRASTRUCTURE-0051 checked against the complete supplied registry and current repository. Distinct upstream records, including repeated questions and broad or damaged extracts, are retained for complete coverage of the source CS category. A repeated record is not a new independent conjecture. Complete mathematical-equivalence and current-literature review has not been performed. Preparation and final review follow contribution prompts C and D, with source gaps disclosed. Structural and rendering checks do not establish a new theorem or open status.
\label{end:20002151}

\section{Literature and context check}
The community-session source was consulted directly. Its entry is a feature or design request, without a complete formal specification. The catalogue title is editorial and is not evidence of an existing theorem. The following finite model is an auxiliary investigation; no PreTeXt implementation or project completion is asserted.
\section{Research attempt: Why stripping a preamble can fail to commute with embedding}
Consider a fragment represented by a pair $(D,B)$ of local definitions and body. Standalone evaluation first installs $D$ and then evaluates $B$. Embedding that deletes $D$ and evaluates $B$ in a host environment is not generally equivalent. Take $D=\{a\mapsto2\}$ and $B=a$, and a host with $a\mapsto3$. Standalone output is 2; stripped embedded output is 3. If the host lacks $a$, the embedded body is undefined instead. Thus removal of the preamble is sound only after proving that its semantic contribution is absent or is supplied identically by the host.

For a finite acyclic macro language one can check the needed condition exactly. Expand the body using the local definitions and list all external macro values actually needed by that expansion. If the host agrees on those values and local names are isolated, structural induction on literals, references, and concatenation proves equal body output in both contexts. Agreement on unused definitions is unnecessary. This is a residual-environment condition, not a version label by itself.

For an overlay, suppose the wrapper cites a base object by an identifier $q$ while a repository has versions $B_1(q)=2$ and $B_2(q)=3$. A bare reference to $q$ cannot determine which body was reviewed. Store a pair $(v,q)$ and resolve it against the selected version. If resolution is a function on that pair, the same pair gives the same base object during standalone and combined assembly. This proves identity consistency under an immutable version map. It does not establish that the wrapper's commentary remains factually appropriate after deliberately selecting a different version.
\section{Verification and adversarial audit}
The argument is an exact analysis of the stated finite model. Its counterexample was checked by direct substitution; no PreTeXt build, browser test, or user study was run. The model's assumptions are explicit and are not asserted to describe every real document.
\paragraph{Audit corrections.} Syntactic preamble removal is not equated with preservation of meaning. A version selector is useful only with defined resolution and retained context, not as a substitute for them.
\paragraph{External computation.} None; symbolic arguments and boundary checks only.
\section{Scope and remaining gaps}
The calculation identifies a concrete failed route or a restricted correctness property. Actual schema integration, transformations, usability, and compatibility with the full original feature request remain unimplemented and unverified. Thus the original agenda remains unresolved, with no claimed mathematical novelty.
\paragraph{Integration review.} Codex (GPT-6; exact serving revision not exposed) reviewed the written arguments and expanded layout and source-reference wrappers. Repository metadata and the production web renderer were checked using Python 3.13 and JavaScript with MathJax 3.2.2. This model-assisted review is not independent human proof verification; no PDF compilation or formal proof checking is asserted.

\section{Completion Estimate}
\noindent\textbf{COMPLETION: 0\%}
The calculations below the preserved source statement investigate a concrete route,
but establish no substantiated new progress toward the entire original target.
Known special cases and failed sufficient conditions are not credited as a solution.
\section{References}
The original source references and catalogue attribution are preserved in the
Sources section above. The following URLs were actually consulted for this attempt;
retention of an additional original reference does not claim fresh consultation of it.
\begin{itemize}
\item \url{https://github.com/PreTeXtBook/community/wiki/Moderated-Problem-Session} (consulted 2026-10-11). 
\end{itemize}
\section{Model and contributor disclosure}
OpenAI gpt-6-astra, requested reasoning configuration ultra, generated this research
attempt on 2026-10-11. The exact serving revision is not exposed. Codex GPT-6
performed orchestration, repository-format review, and a model-assisted mathematical
audit of the written arguments; its exact serving revision
is also not exposed. The mathematical route, boundary checks and final adversarial
self-audit followed contribution prompts A then B in
\texttt{docs/contribution\_prompts.md}, after reading \texttt{CONTRIBUTING.md}.
Python 3.14.2 assembled the standalone documents; any mathematical computation is
specified in Verification. No theorem prover was used. No human mathematical
authorship, substantial human proof correction, or independent human proof checking
is asserted. Model-assisted review, source retrieval, and structural checks do not
establish correctness or novelty. Full private reasoning transcripts are unavailable
as reproducibility artifacts; the complete arguments, inputs and outcomes actually
used are stated here or in the companion batch audit.
\section{Final status}
Unresolved as a new research contribution. This is a concrete route investigation,
not a claim of new complete resolution or a current-status certification. A known
answer reproduced for context does not become a new contribution. The catalogue's
source-status label is preserved; this attempt status does not assert that a
source-labelled solved question is open.
\end{document}
